The modular level is the conductor of these realizations by the
published result recalled below. Li's Fricke relation and Shimura's
Jacobian construction, as in the baseline modularity argument, also
place $r_5$ in the Tate module of $J_1(N)$.

\paragraph{The primitive level equals the local conductor.}
Let $f\in S_2(\Gamma_0(N),\eta)$ be the primitive newform obtained
after absorbing the finite character twist, and let $U_{f,\ell}$
be its two-dimensional Tate realization for the chosen coefficient
embedding. Here $\eta$ is the nebentypus of $f$; it is not identified
with the controlled character $\chi=b^2$ in (LR2).
Carayol's published conductor result for normalized primitive
newforms gives
\[
 a_p(U_{f,\ell})=v_p(N),\qquad p\ne\ell.
\tag{CL1}
\]
Indeed, \cite[\S2.1, p.~36]{carayol1987} permits arbitrary
nebentypus and coefficient number fields, and weight two.
Section~2.4 on p.~37 states the conductor assertion also at primes
dividing $N$; the conclusion following the full Theorem~4 in
\S4.1 on p.~42 proves that assertion. We use the full theorem,
not its weaker version with an extraordinary-cuspidal exclusion.
Thus the assertion includes dyadic primes and nonsquarefree level.
The cohomological realization in that source is dual to the Tate
realization; duality preserves the Artin conductor. The unramified
norm twist in its Hecke-correspondence convention likewise changes
no conductor exponent.

For each rational prime $p$, choose an auxiliary prime $\ell\ne p$.
The preceding compatible-family comparison identifies the actual
$r_\ell$ constructed from $A$, $a_\sigma$, and $b$ with
$U_{f,\ell}$ for this same final newform and coefficient embedding.
Restriction to $G_{\mathbb Q_p}$ preserves this isomorphism, hence
preserves inertia-invariant dimensions and the Swan conductor.
Consequently
\[
 a_p(r_\ell)=a_p(U_{f,\ell})=v_p(N),\qquad p\ne\ell.
\tag{CL2}
\]
The conductor includes the tame term and Swan term, or equivalently
the full Weil--Deligne conductor with its monodromy contribution.
No semisimplification on inertia is taken. Applying (CL2) with an
auxiliary prime different from $p$ controls every prime divisor of
$N$; no local compatibility theorem at $p=\ell$ is required.

